% D.93 integrated after full root, FGH and PBC proof/code audits of the original theorem and all-depth Smith addendum.
\subsection{Unbounded finite compact kernels on actual ordinary fibres}

\begin{statement}{D.93}{Proved}{All finite compact kernels in an ordinary-point family}
For every integer \(r\ge1\), let
\[
 t_r=2^{5^r},\qquad
 s_r=\bigl(\kappa(t_r),\kappa(1-t_r)\bigr),
 \qquad n_r=5^r.                                      \tag{D.93.1}
\]
This is the abelian image of the actual rational-point section
of the tripod at \(t_r\). Retain the compact groups
\(\mathcal S_s\) and \(\mathcal K_s\) of D.63.
On the free \(\mathbb Z_5\)-module with basis
\(e_j\), \(j\in\mathbb Z/n_r\), define the two literal
permutation matrices
\[
 T_re_j=e_{j+1},\qquad
 \Pi_re_j=e_{2j},
 \qquad
 M_r=2(I+T_r)^3\Pi_r.                                 \tag{D.93.2}
\]
Then the entire separated compact kernel has the exact finite
presentation
\[
 \begin{gathered}
 \mathcal K_{s_r}\simeq\mathcal S_{s_r}
 \simeq\operatorname{coker}_{\mathbb Z_5}(M_r-I)
 \simeq
 \left(\bigoplus_{j=1}^{r-1}\mathbb Z/5^j\right)
       \oplus\mathbb Z/5^{r+1},\\
 \#\mathcal S_{s_r}=5^{\,1+r(r+1)/2}.
 \end{gathered}                                      \tag{D.93.3}
\]
In particular these finite orders are unbounded on actual
ordinary rational-point fibres. The first two integral Smith
calculations give
\[
 \mathcal K_{s_1}\simeq\mathbb Z/25,\qquad
 \mathcal K_{s_2}\simeq
        \mathbb Z/5\ \oplus\ \mathbb Z/125.             \tag{D.93.4}
\]
Equation (D.93.2) is a fully explicit finite integral matrix
for every \(r\). The all-depth elementary divisors follow from
the integral norm-extension proof below; the two small Smith
calculations are independent checks.

Every indicated metabelian fibre is nonempty and has translation
group \(\prod_{\mathbb N}\mathbb Z_5\), with its compact product
topology. Exactly one ordinary rational point has its abelian
image, namely \(t_r\); no cusp has that abelian image. Thus every
other class in this metabelian fibre is neither a rational-point
class nor a cusp class. This is a statement about actual
geometrically pro-five metabelian sections. It does not classify
their higher-derived lifts or general abelian-pair existence.
\end{statement}

\begin{proof}
\paragraph{The global lattice and the exact local depths.}
Write \(B_r=2^{5^r}-1>0\). Since the sign has zero pro-five
Kummer class, the pair in (D.93.1) is
\[
        s_r=\bigl(5^r\kappa(2),\kappa(B_r)\bigr).
                                                               \tag{D.93.5}
\]
The integer \(B_r\) is odd and prime to five. It is not a
fifth power: put \(h=2^{5^{r-1}}\), so that
\((h-1)^5<h^5-1<h^5\). Consequently at least one prime
valuation of \(B_r\) is a unit in \(\mathbb Z_5\).
The first coordinate is supported only at two, and two does
not divide \(B_r\). Hence the actual closed away-five lattice is
\[
        L=5^r\mathbb Z_5(1,0)\oplus\mathbb Z_5(0,1),
             \qquad (a_5,b_5)=(0,0).                  \tag{D.93.6}
\]
This uses the complete valuation lattice, not only its reduction
modulo five. That reduction has rank one, so the earlier
residual-rank-two results do not give (D.93.3).

For a rational five-adic unit \(x\), use the integral coefficient
\[
                   \lambda_x=\frac{\log(x^4)}{4\log6}.
\]
It is the coefficient of \(\kappa(6)\) in the local unit class.
The denominator has valuation one. Direct reduction modulo
twenty-five gives
\[
 \lambda_2\equiv2\pmod5,\qquad
 2^{5^r}\equiv7\pmod{25},\qquad
 1-2^{5^r}\equiv-6\pmod{25},
 \qquad
 \lambda_{1-2^{5^r}}\equiv1\pmod5.                    \tag{D.93.7}
\]
Indeed \(5^r\equiv5\pmod{20}\), while two has order twenty
modulo twenty-five. Also
\(16-1\equiv15\pmod{25}\) and \(6^4-1\equiv20\pmod{25}\),
which give the two logarithmic congruences.
Put
\[
 \alpha=\lambda_2,\qquad
 \beta=\lambda_{1-2^{5^r}},\qquad
 \mathbf z=(5^r\alpha,\beta),\qquad
 \gamma=-\alpha/\beta\in\mathbb Z_5^\times .
                                                               \tag{D.93.8}
\]
These are the actual local coefficients in D.87.
The three global depths of \(a,b,a-b\) are respectively
\((r,0,0)\), and their three local depths are also \((r,0,0)\).
The first assertion follows from (D.93.6); the second follows
from (D.93.7), including the unit \(\beta-5^r\alpha\).
D.77 therefore gives
\[
                       \mathcal K_{s_r}\simeq
                       \mathcal S_{s_r}.              \tag{D.93.9}
\]
There is no discarded local free factor in this equality.

\paragraph{The actual image comparison and its natural map.}
Use \(\Lambda=\mathbb Z_5[[u,v]]\), \(t_1=1+u\),
\(t_2=1+v\), \(w=t_1t_2-1\), and \(P=uvw\).
Let \(I_0\) be the actual untwisted Alexander operator image,
including its cyclotomic coordinate \(\chi\), and let
\(N_0=\ker(\chi:I_0\to\mathbb Z_5^\times)\).
There is no valuation graph here, since (D.93.6) has zero
five-coordinate. D.87 supplies the actual common-complement
images
\[
              I_s=L\rtimes I_0,\qquad
              I_5=\mathbb Z_5\mathbf z\rtimes I_0 .
                                                               \tag{D.93.10}
\]
It also identifies the finite-localization kernel with its
natural coefficient map. In particular one finite coefficient
level suffices for all the local trivializations used in this
identification; no unrestricted locally trivial group is being
substituted for \(\mathcal S_s\).

Set \(q=(1+u)^{5^r}-1\). The local coefficient quotient and
translation ideal are
\[
 R=\Lambda/(t_1^{5^r\alpha}t_2^\beta-1),
                    \qquad J=(q,v)R.                 \tag{D.93.11}
\]
Since \(\beta\) is a unit, the first relation eliminates \(t_2\)
integrally. Thus, as actual topological rings,
\[
 R=\mathbb Z_5[[u]],\qquad
 t_2=t_1^{5^r\gamma},\qquad
 v=(1+q)^\gamma-1,\qquad J=qR.                         \tag{D.93.12}
\]
The series \(v/q\) has unit constant coefficient \(\gamma\).
Likewise
\[
 w=t_1^{1+5^r\gamma}-1
\]
is \(u\) times a unit, since \(1+5^r\gamma\) is a unit.
Consequently the ideal generated by the actual \(P\) in \(R\) is
\[
                              PR=u^2qR.               \tag{D.93.13}
\]
No rationalization or splitting into cyclotomic characters
has been used in these integral identities.

For every \(g\in N_0\), the exact beta boundary identities
of D.39 give \(\mathcal B_g-1\in(P)\). The primitive actual
local cubic element in D.77 has
\(\mathcal B_\delta-1=P E_\delta\), with \(E_\delta\) a unit.
Hence the ideal of all these actual operator differences is
exactly \((P)\), on \(R\) and on the invariant ideal \(J\).
Ordinary duality applied to D.87.4 therefore gives
\[
 \mathcal S_s=
 \ker\left\{
       (J/PJ)(-1)_\Gamma\longrightarrow
       (R/PR)(-1)_\Gamma\right\},
                    \qquad \Gamma=\mathbb Z_5^\times .
                                                               \tag{D.93.14}
\]
The map is still induced by the actual inclusion \(J\subset R\).
The twist \((-1)\) records precisely the Cartier-to-ordinary
duality convention. These are compact coinvariants, with closed
images, and (D.93.14) is a map of compact modules.

\paragraph{Retain the generator before taking the two low jets.}
Scalar division by the retained generator \(q\) identifies
\(J/PJ\) with
\[
                         T=R/(u^2q).
\]
On this scalar module the actual Cartier-untwisted action is
\[
        \chi(g)^{-1}S_{\chi(g)}(t_1^{5^r})\rho_0(g),
          \qquad S_c(t)=\frac{t^c-1}{t-1},             \tag{D.93.15}
\]
where \(\rho_0\) includes its substitution \(t_i\mapsto
t_i^{\chi(g)}\). The natural map in (D.93.14) is multiplication
by \(q\), not the identity on two isomorphic scalar quotients.

The quotient \(T/u^2T=R/(u^2)\) is free of rank two over
\(\mathbb Z_5\). Its successive actual characters are
\(\chi^2,\chi^3\): the Alexander multiplier has augmentation
\(\chi^2\), the retained \(q\) contributes \(\chi\), the
Cartier untwist removes one factor \(\chi\), and the linear
substitution supplies the next factor. Both characters are
nontrivial on the order-four tame subgroup. Exact averaging
by that subgroup shows that every continuous
\(\Gamma\)-homology group of this two-jet quotient is zero.
The compact exact homology sequence consequently gives an
isomorphism
\[
                 (u^2T)_\Gamma\xrightarrow{\ \sim\ }T_\Gamma.
                                                               \tag{D.93.16}
\]
This proves both surjectivity and injectivity; no claim that a
coinvariant functor is generally left exact is needed.

Since \(R\) is a domain, scalar division by \(u^2\) gives
\[
                u^2T\simeq R/(q)
                    \simeq\mathbb Z_5[C_{5^r}].        \tag{D.93.17}
\]
On this quotient \(t_1^{5^r}=t_2=1\).
The exact untwisted Alexander action is
\[
 \rho_0(g)F=
 S_\chi(t_1)S_\chi(t_2)\mathcal B_g(t_1,t_2)
                 F(t_1^\chi,t_2^\chi).
\]
The boundary identity \(\mathcal B_g(t,1)=1\) is exact for
every multiplier. Combining it with (D.93.15) and the
substitution on the retained \(u^2\) gives on (D.93.17)
\[
                       F(t)\longmapsto
                  \chi\,S_\chi(t)^3 F(t^\chi).         \tag{D.93.18}
\]
Thus neither an unspecified Frobenius constant nor an
approximate weight action remains.

Multiplication by \(q\) sends the actual \(u^2\)-layer to
zero in \(R/PR\), since \(u^2q=0\) there. By (D.93.16)
the entire natural map in (D.93.14) is zero.
It follows that
\[
             \mathcal S_{s_r}\simeq
       \bigl(\mathbb Z_5[C_{5^r}],
                  \chi S_\chi(t)^3\sigma_\chi\bigr)_\Gamma .
                                                               \tag{D.93.19}
\]
The element two topologically generates \(\mathbb Z_5^\times\):
its order modulo five is four and
\(v_5(2^4-1)=1\). Its actual action in (D.93.18) is precisely
the matrix \(M_r\) of (D.93.2). The image of \(M_r-I\)
on this finite free compact module is closed. Density of the
cyclic subgroup generated by two therefore identifies its
cokernel with the full compact \(\Gamma\)-coinvariants.
This proves the finite presentation in (D.93.3).

\paragraph{The all-depth determinant.}
Compute the determinant after adjoining the \(5^r\)-th roots
of unity in characteristic zero; this does not change it.
The evaluation at \(t=1\) has scalar \(16\).
For each \(1\le j\le r\), multiplication of the exponents by
two is one orbit on the primitive \(5^j\)-th roots, of length
\(\ell_j=4\cdot5^{j-1}\). On that orbit the product of the
weights in (D.93.18) is
\[
 2^{\ell_j}\prod_{k=0}^{\ell_j-1}
       (1+\zeta^{2^k})^3
 =2^{\ell_j},
 \qquad
 1+\zeta^{2^k}
     =\frac{\zeta^{2^{k+1}}-1}{\zeta^{2^k}-1}.          \tag{D.93.20}
\]
The product telescopes because the orbit closes and none
of these primitive roots is one. For a weighted permutation
cycle the determinant of identity minus the cycle is one
minus its weight product. Consequently
\[
 \det(I-M_r)
      =(1-16)\prod_{j=1}^r
                 (1-2^{\,4\cdot5^{j-1}}).             \tag{D.93.21}
\]
All factors are nonzero. The elementary lifting-the-exponent
calculation gives
\[
 v_5(16-1)=1,\qquad
 v_5(2^{\,4\cdot5^{j-1}}-1)=j.
 \]
For example this follows inductively by expanding
\((1+5c)^5-1\), which increases the valuation by exactly
one when \(c\) is a unit. Thus the cokernel is finite of the
order in (D.93.3), with its finite compact topology.
The exact literal integral Smith forms for \(r=1,2\) are
executed in the accompanying certificate and give (D.93.4).
Those two executions are not used to prove (D.93.21).

\paragraph{An exact extension by one order-five factor.}
Put \(n=5^r\), \(R=\mathbb Z_5[t]/(t^n-1)\), \(u=t-1\),
and \(I=uR\). Let \(\sigma_\chi(t)=t^\chi\).
On the same underlying free module \(R\), distinguish
\[
 \mathscr A:\ \chi\sigma_\chi,\qquad
 \mathscr B:\ \chi S_\chi(t)^3\sigma_\chi,
       \qquad S_\chi(t)=\frac{t^\chi-1}{t-1}.
\]
The module \(\mathscr B_\Gamma\), with
\(\Gamma=\mathbb Z_5^\times\), is exactly the compact kernel
in D.93. Write
\[
                  \mathcal N=1+t+\cdots+t^{n-1}.
\]
Multiplication by \(u^3\) gives the exact equivariant sequence
\[
 0\longrightarrow \mathbb Z_5\mathcal N(\chi^4)
   \longrightarrow \mathscr B
   \xrightarrow{u^3} I^3\mathscr A
   \longrightarrow0.                                \tag{D.93.22}
\]
Indeed the equivariance is
\(u^3 S_\chi(t)^3=\sigma_\chi(u^3)\).
The kernel of multiplication by \(u^3\) in \(R\) is precisely
\(\mathbb Z_5\mathcal N\): over a characteristic-zero splitting
field it is supported only at \(t=1\), hence is killed by \(u\);
the coefficient equations for \(u f=0\) then say that all
coefficients of \(f\) are equal. This also proves the integral
kernel assertion. The action on this generator is \(\chi^4\).

The quotient \(\mathscr A/I^3\mathscr A\) has the three
successive tame characters \(\chi,\chi^2,\chi^3\), on its
augmentation filtration. These assertions include any
five-torsion in its two positive-degree pieces. All three
characters are nontrivial on the order-four subgroup of
\(\Gamma\). Exact tame averaging therefore kills every
\(\Gamma\)-homology group of the quotient. Consequently
\[
 (I^3\mathscr A)_\Gamma\simeq\mathscr A_\Gamma,
 \qquad H_1(\Gamma,I^3\mathscr A)=H_1(\Gamma,\mathscr A)=0.
                                                               \tag{D.93.23}
\]
For the last equality, the completed procyclic resolution
after tame averaging identifies \(H_1\) with invariants.
The action of two on \(\mathscr A\) is twice a finite
permutation, so its fixed module is zero. All modules here
are finite free compact \(\mathbb Z_5\)-modules before the
displayed finite quotients; the compact exact resolution is
being used with its actual topology.

Decompose the basis \(t^a\) of \(\mathscr A\) into its orbits
under \(a\mapsto2a\) in \(\mathbb Z/n\).
The zero orbit has weight two and contributes nothing.
For each \(1\le j\le r\), the orbit of
\(t^{5^{r-j}}\) has length \(\ell_j=4\cdot5^{j-1}\);
its weighted-cycle cokernel is
\(\mathbb Z_5/(2^{\ell_j}-1)\simeq\mathbb Z/5^j\).
The exact homology sequence of (D.93.22) now gives
\[
 0\longrightarrow\mathbb Z/5\,[\mathcal N]
 \longrightarrow\mathscr B_\Gamma
 \longrightarrow\bigoplus_{j=1}^r\mathbb Z/5^j
 \longrightarrow0.                                  \tag{D.93.24}
\]
In particular the norm class has exact order five. It is
not discarded by a rational calculation.

\paragraph{An explicit lift of the largest orbit generator.}
Let \(\sigma=\sigma_2\), and write
\[
 A=2\sigma,\qquad B=2(1+t)^3\sigma
\]
for the actual operators on \(\mathscr A,\mathscr B\).
The integral-at-five polynomials
\[
 G=-1-\frac{u}{3}+\frac{2u^2}{21},
       \qquad f=\frac4{21}(t+3)
\]
satisfy the literal polynomial identity
\[
                 u^3 f=t+(A-1)G.                     \tag{D.93.25}
\]
Thus \([f]\) maps in (D.93.24) to the largest orbit
generator \([t]\), up to the stated natural identification.

Put
\[
 \ell=4\cdot5^{r-1},\qquad
 c=2^\ell-1,\qquad
 Q=\sum_{i=0}^{\ell-1}2^i t^{2^i},\qquad H=Q+cG .
                                                               \tag{D.93.26}
\]
Here \(v_5(c)=r\) and \(c/n\equiv3\pmod5\).
The weighted orbit sum gives
\[
                       (A-1)H=c\,u^3f
                          \quad\hbox{in }R.           \tag{D.93.27}
\]
Indeed \(A^\ell t=2^\ell t^{2^\ell}=(1+c)t\) in
\(R\), since \(n\mid c\). This argument is in the actual
finite group ring; it does not set \(t^{2^\ell}=t\)
in the unrestricted polynomial ring.

We give an integral lift of \(H\) through multiplication
by \(u^3\). Write \(h_i\) for its coefficient of \(u^i\)
in the unrestricted polynomial ring. Its constant
coefficient is zero. For
\[
        S_k=\sum_{i=0}^{\ell-1}2^{ki}
               =\frac{(1+c)^k-1}{2^k-1},
\]
its first three coefficients are
\[
 \begin{aligned}
 h_1&=\frac{c+c^2}{3},\\
 h_2&=-\frac c{42}+\frac{c^2}{21}+\frac{c^3}{14},\\
 h_3&=\frac{S_4-3S_3+2S_2}{6}.
 \end{aligned}                                       \tag{D.93.28}
\]
The first two lie in \(n\mathbb Z_5\). Therefore
\[
 W_0=h_1/n,\qquad
 W_1=\bigl(h_2-\tbinom n2W_0\bigr)/n
                       \quad\hbox{belong to }\mathbb Z_5,
\]
and there is an actual polynomial \(Z\in\mathbb Z_5[u]\)
with
\[
 H-(t^n-1)(W_0+W_1u)=u^3Z.                            \tag{D.93.29}
\]
Since \(v_5(\binom n2)=v_5(\binom n3)=r\), its
constant coefficient \(z_0=Z|_{u=0}\) satisfies
\[
                         z_0\equiv h_3\pmod{5^r}.      \tag{D.93.30}
\]

\paragraph{The norm extension is nonzero on the largest factor.}
Equations (D.93.27)--(D.93.29), and the integral kernel
calculation in (D.93.22), give
\[
             c f-(B-1)Z=k\mathcal N
                   \quad\hbox{in }R,\quad k\in\mathbb Z_5.
                                                               \tag{D.93.31}
\]
Evaluation at \(t=1\) yields the exact equality
\[
                 nk=c\,\frac{16}{21}-15z_0.            \tag{D.93.32}
\]
We compute \(k\) modulo five. By (D.93.30), replacing
\(z_0\) by \(h_3\) changes \(15z_0/n\) by a multiple
of five. The terms \(S_2,S_3\) in (D.93.28) have
valuation at least \(r\), so they disappear after this
operation modulo five. The \(S_4\) term gives exactly
\[
 \frac{15h_3}{n}
   \equiv\frac{(1+c)^4-1}{6n}
   \equiv\frac{2c}{3n}\pmod5.
\]
All expressions compared here are integral at five.
Consequently
\[
       k\equiv\frac cn\left(\frac{16}{21}-\frac23\right)
        =\frac{2c}{21n}\equiv1\pmod5.                 \tag{D.93.33}
\]
Thus \(c[f]\) is a nonzero norm class in (D.93.24).
Since \(v_5(c)=r\), the lift \([f]\) has exact order
\(5^{r+1}\). This proves the required nonsplitting over
the largest cyclic quotient.

\paragraph{Splitting every smaller factor.}
Rescale this lift by a five-adic unit to obtain \(x_r\)
with \(5^r x_r=[\mathcal N]\). Choose arbitrary lifts
\(x_j\) of the smaller orbit generators, \(1\le j<r\).
For some \(a_j\in\mathbb Z/5\), they satisfy
\(5^jx_j=a_j[\mathcal N]\).
Choose integral representatives of \(a_j\) and replace
\[
                     x_j\ \longmapsto\
                     x_j-a_j5^{r-j}x_r .
 \]
The new lift has order exactly \(5^j\). On the quotient
in (D.93.24), these replacements are invertible triangular
basis changes: the \(r\)-th generator is kept, and each
smaller one is changed only by a permissible multiple of
that generator. The modified lifts and \(x_r\) generate
the whole module. Their product of orders equals the
order supplied by (D.93.24). They therefore give precisely
the direct-sum decomposition (D.93.3).

This closes the all-depth elementary-divisor calculation.
The literal symbolic identities and moment checks are retained
in \path{code/de_ordinary_kernel_smith.py}.
The small integral matrices are retained in
\path{code/de_ordinary_finite_kernel.py}.
The extension argument above supplies the proof for every \(r\).

\paragraph{The actual section fibres and their ordinary points.}
The rational point \(t_r\) supplies a full section, hence an
actual metabelian class above \(s_r\). The pair has finite
prime support. D.74 therefore identifies the translation
group of its metabelian fibre with
\(\prod_{\mathbb N}\mathbb Z_5\).
The two Kummer coordinates are linearly independent over
\(\mathbb Q_5\), by their disjoint nonempty prime supports.
They lie on none of the three completed cusp lines in D.35.

For a nonzero rational number \(x\), its pro-five Kummer
class vanishes exactly when \(x=\pm1\); this follows directly
from its prime valuations. If an ordinary rational point
\(z\) has abelian pair \(s_r\), it follows that
\[
                  z=\pm t_r,\qquad
                  1-z=\pm(1-t_r).
 \]
The choice \(z=-t_r\) makes the second equality impossible
in either sign. Hence \(z=t_r\), as claimed. These statements
classify the ordinary point images in this fibre, not the
full section classes of the geometric fundamental group.
\end{proof}
